\documentclass[11pt,a4paper]{article}
\usepackage[T1]{fontenc}
\usepackage{lmodern}
\usepackage[margin=22mm,headheight=14pt]{geometry}
\usepackage{amsmath,amssymb,graphicx,booktabs,tabularx,array}
\usepackage{xcolor,listings,enumitem,microtype,fancyhdr}
\usepackage{tikz}
\usetikzlibrary{arrows.meta,positioning}
\usepackage[font=small,labelfont=bf]{caption}
\captionsetup{hypcap=false}
\usepackage[hidelinks]{hyperref}
\hypersetup{pdftitle={MIT VNAV Individual Deliverables},
pdfauthor={yawangjiang}}
\definecolor{navy}{HTML}{19394F}
\definecolor{pale}{HTML}{F0F4F7}
\setlength{\parindent}{0pt}
\setlength{\parskip}{5pt}
\setlength{\emergencystretch}{2em}
\setlist[itemize]{leftmargin=1.5em,itemsep=3pt,topsep=3pt}
\setlist[enumerate]{leftmargin=1.7em,itemsep=3pt,topsep=3pt}
\pagestyle{fancy}
\fancyhf{}
\fancyhead[L]{\small\color{navy}MIT VNAV - Individual}
\fancyhead[R]{\small\color{navy}yawangjiang}
\fancyfoot[C]{\small\thepage}
\renewcommand{\headrulewidth}{0.3pt}
\lstset{basicstyle=\ttfamily\small,columns=fullflexible,keepspaces=true,
breaklines=true,breakatwhitespace=false,showstringspaces=false,
backgroundcolor=\color{pale},frame=single,rulecolor=\color{pale},
xleftmargin=5pt,xrightmargin=5pt,framesep=5pt,aboveskip=7pt,belowskip=7pt}
\newcommand{\code}[1]{\texttt{#1}}
\newcommand{\R}{\mathbb{R}}
\newcommand{\ex}{\mathbf e_x}
\newcommand{\ev}{\mathbf e_v}
\newcommand{\eR}{\mathbf e_R}
\newcommand{\eO}{\mathbf e_\Omega}
\newcommand{\eThree}{\mathbf e_3}
\newcommand{\Fdes}{\mathbf F_{\mathrm{des}}}
\newcommand{\Fdrag}{\mathbf F_{\mathrm{drag}}}
\newcommand{\repo}{https://github.com/LittleHorseyaya/vnav}

\hypersetup{pdftitle={MIT VNAV LAB4 - Individual Deliverables 1 and 2},pdfauthor={yawangjiang}}
\begin{document}
\begin{center}
{\LARGE\bfseries\color{navy}LAB4: Trajectory Optimization}\par
\vspace{6pt}
{\large Individual Deliverables 1 and 2}\par
\vspace{6pt}
Submission account: \code{yawangjiang}\par
October 31, 2026
\end{center}

\section{Deliverable 1: Single-Segment Optimization}
The problem is to minimize $J=\int_0^1[P'(t)]^2\,dt$, with $P(0)=0$ and
$P(1)=1$. Write $P(t)=\sum_{i=1}^n p_i t^i$, so that $p_0=0$ is eliminated.
Then
\begin{equation}
J=\mathbf p^\top Q_n\mathbf p,\qquad
(Q_n)_{ij}=\frac{ij}{i+j-1},\quad 1\leq i,j\leq n.
\end{equation}
The objective uses no factor $1/2$; a solver using
$\frac12\mathbf p^\top H\mathbf p$ should receive $H=2Q_n$.

\subsection{1. Degree One}
For $P(t)=p_1t$, the final-position constraint gives $p_1=1$. Thus
\begin{equation}
\boxed{P(t)=t,\qquad J=1.}
\end{equation}

\subsection{2. Degree Two}
For $\mathbf p=[p_1,p_2]^\top$,
\begin{equation}
Q_2=\begin{bmatrix}1&1\\1&4/3\end{bmatrix},\qquad
A_2=\begin{bmatrix}1&1\end{bmatrix},\qquad b=1.
\end{equation}
In particular, $J=p_1^2+2p_1p_2+\frac43p_2^2$. Substituting
$p_1=1-p_2$ yields $J=1+\frac13p_2^2$. Therefore
\begin{equation}
\boxed{\mathbf p^*=[1,0]^\top,\qquad P(t)=t,\qquad J^*=1.}
\end{equation}
Allowing the quadratic term does not reduce the optimum.

\subsection{3. Degree Three and the Euler-Lagrange Result}
For $\mathbf p=[p_1,p_2,p_3]^\top$,
\begin{equation}
Q_3=\begin{bmatrix}1&1&1\\1&4/3&3/2\\1&3/2&9/5\end{bmatrix},\qquad
A_3=\begin{bmatrix}1&1&1\end{bmatrix},\qquad b=1.
\end{equation}
The minimizer is $\mathbf p^*=[1,0,0]^\top$, again giving $P(t)=t$ and $J^*=1$.
Indeed, Cauchy-Schwarz gives
\begin{equation}
1=\left(\int_0^1P'(t)\,dt\right)^2\leq\int_0^1[P'(t)]^2\,dt.
\end{equation}
Equality holds for the constant derivative $P'=1$. This agrees with the
Euler-Lagrange equation $P''=0$ and the two position boundary conditions.

\clearpage
\subsection{4. Add the Terminal-Velocity Constraint \texorpdfstring{$P'(1)=-2$}{P'(1)=-2}}
The coefficient relation for this additional condition is
$\sum_{i=1}^n i p_i=-2$. For each degree, solve
\begin{equation}
\min_{\mathbf p}\mathbf p^\top Q_n\mathbf p
\quad\text{subject to}\quad A\mathbf p=\mathbf b.
\end{equation}
One direct method is the equality-constrained KKT system
\begin{equation}
\begin{bmatrix}2Q_n&A^\top\\A&0\end{bmatrix}
\begin{bmatrix}\mathbf p\\\boldsymbol\lambda\end{bmatrix}
=\begin{bmatrix}\mathbf0\\\mathbf b\end{bmatrix}.
\end{equation}

\subsubsection*{4(a). Degree Two}
Use the same $Q_2$ and
\begin{equation}
A=\begin{bmatrix}1&1\\1&2\end{bmatrix},\qquad
\mathbf b=\begin{bmatrix}1\\-2\end{bmatrix}.
\end{equation}
The constraints give $p_1+p_2=1$ and $p_1+2p_2=-2$, hence
\begin{equation}
\boxed{P(t)=4t-3t^2,\qquad J=\int_0^1(4-6t)^2\,dt=4.}
\end{equation}
There is only one feasible quadratic coefficient vector.

\subsubsection*{4(b). Degree Three}
Use $Q_3$ and
\begin{equation}
A=\begin{bmatrix}1&1&1\\1&2&3\end{bmatrix},\qquad
\mathbf b=\begin{bmatrix}1\\-2\end{bmatrix}.
\end{equation}
Solving the KKT equations gives
\begin{equation}
\mathbf p^*=\begin{bmatrix}1/4\\9/2\\-15/4\end{bmatrix},\qquad
\boxed{P(t)=\frac14t+\frac92t^2-\frac{15}{4}t^3,\qquad J^*=\frac{17}{8}.}
\end{equation}
The boundary values satisfy $P(0)=0$, $P(1)=1$, and
$P'(1)=\frac14+9-\frac{45}{4}=-2$. Its cost is $17/8=2.125<4$, so the
additional cubic freedom improves the constrained quadratic solution.

\subsection*{Summary of Deliverable 1}
\begin{center}\small
\begin{tabular}{@{}lcc@{}}
\toprule
Allowed degree and constraints & Optimal polynomial & Cost \\
\midrule
1, 2, or 3; endpoint positions only & $t$ & $1$ \\
2; also $P'(1)=-2$ & $4t-3t^2$ & $4$ \\
3; also $P'(1)=-2$ & $\frac14t+\frac92t^2-\frac{15}{4}t^3$ & $17/8$ \\
\bottomrule
\end{tabular}
\end{center}

\clearpage
\section{Deliverable 2: Multi-Segment Optimization}
Assume fixed, positive segment durations, specified waypoint positions, fixed
derivatives through order $r-1$ at the two outer endpoints, and free derivatives
of orders $1,\ldots,r-1$ at the internal waypoints. Counts below are per scalar
coordinate. Independent spatial coordinates repeat the same calculation.

\subsection{1. Minimum Snap: \texorpdfstring{$r=4$ and $k$}{r=4 and k} Segments}
For $J=\sum_{i=1}^k\int[P_i^{(4)}(t)]^2\,dt$, the Euler-Lagrange equation is
$P_i^{(8)}=0$. Each segment therefore has degree at most 7 and 8 coefficients,
giving $8k$ unknown coefficients.

Count the equations using the following disjoint categories:
\begin{center}\small
\begin{tabularx}{\linewidth}{@{}>{\raggedright\arraybackslash}p{35mm}c>{\raggedright\arraybackslash}X@{}}
\toprule
Category & Equations & Meaning \\
\midrule
Waypoint positions & $2k$ & Two endpoint position values for each segment, representing $k+1$ distinct geometric waypoints. \\
Fixed nonzero-order derivatives & $6$ & Velocity, acceleration, and jerk at the overall start and finish. \\
Free internal derivatives & $6(k-1)$ & Each of the three free derivatives at every internal waypoint yields continuity and a natural optimality equation. \\
\bottomrule
\end{tabularx}
\end{center}
At an internal waypoint, freedom in derivatives $j=1,2,3$ gives, respectively,
continuity of derivative orders $(1,6)$, $(2,5)$, and $(3,4)$. The second member
of each pair is the natural stationarity condition obtained by integration by
parts. Thus
\begin{equation}
2k+6+6(k-1)=8k.
\end{equation}
Equivalently, classify the internal positions separately: $2(k-1)$ internal
position equations, 8 fixed endpoint conditions of orders 0 through 3, and
$6(k-1)$ free-derivative equations. Endpoint positions must not be counted twice.

\subsection{2. General Derivative Order \texorpdfstring{$r$}{r}}
The Euler-Lagrange equation is $P_i^{(2r)}=0$, so each segment has degree at most
$2r-1$ and $2r$ coefficients, for $2rk$ unknowns. Internal free derivative $j$
gives continuity of order $j$ and the natural matching condition of order
$2r-1-j$. Consequently
\begin{equation}
\boxed{2k+2(r-1)+2(r-1)(k-1)=2rk.}
\end{equation}
The three terms count waypoint positions, fixed outer derivatives of orders
$1,\ldots,r-1$, and internal free-derivative continuity/optimality equations.
The equivalent endpoint-inclusive count is
$2(k-1)+2r+2(r-1)(k-1)=2rk$.

For $r=4,k=3$, this gives $6+6+12=24$ equations for 24 coefficients. A free
derivative is not an extra prescribed value: its equations follow from
continuity and optimality. The highest derivative $2r-1$ may jump at a fixed
internal position waypoint.

\subsection*{Source}
MIT VNAV LAB4 exercise statement:
\url{https://vnav.mit.edu/labs_2023/lab4/exercises.html}.
\end{document}
